\scanpage{75}
\mspage{68}
% archive sheet number: 74
\diagram{
\begin{tikzcd}[ampersand replacement=\&,column sep=large,row sep=large]
M(I')_cW(I')_c^{-1}\arrow[r,dashed,"\simeq"]\&\deleted{\unsure{M(I')}}\mathrm{HO}_W(\unsure{I'|I})\&M(I')_f\arrow[l,dashed,"\simeq"']\arrow[d,"u_*^W"]\\
M(I)_cW(I)_c^{-1}\arrow[u,"u_W^{\unsure{*|\deleted{c}*}}"]\&\mathrm{HO}^{\deleted{\unsure{f|\circ}}}_W(I)\arrow[l,dashed,"\simeq"']\&M(I)_f\arrow[l,dashed,"\simeq"']
\end{tikzcd}
}
Et on veut montrer que $u_W^*$ et $u_*^W$ sont adjoints, de façon précise trouver un isom. bifonctoriel, pour \uline{$F$ dans $M(I)_c$} et \uline{$G'$ dans $M(I')_f$} \deleted{\unsure{iso}}
\[
\operatorname{Hom}_W(u^*F,\unsure{C'|G'})\xrightarrow{\sim}\operatorname{Hom}_W(F,u_*G')
\tag*{($\alpha$)}
\]
(où $\operatorname{Hom}_W$ désigne les $\operatorname{Hom}$ de $\mathrm{HO}_W(I')$ resp $\mathrm{HO}_W(I)$). Comme $F$ est dans $M(I)_c$, $u^*F$ est dans $M(I')_c$, et comme $G'$ est dans $M(I')_f$, $u_*(G')$ est dans $M(I)_f$. Il s'ensuit que les deux mbres de $(\alpha)$ sont des quotients des deux mbres dans la bijection canonique
\[
\operatorname{Hom}_{M(I')}(u^*F,G')\overset{\alpha}{\xrightarrow{\sim}}\operatorname{Hom}_{M(I)}(F,u_*G').
\tag*{($\alpha\alpha$)}
\]
Il faut simplement montrer que cette dernière transforme \unsure{la relation|les relations} d'équivalence pertinente $R'$ du premier membre,
